% Generated by task tex-opening from dossiers/gt-opening-compile.md (sections E, F and G).
% With the corrections of dossiers/verify-gt-opening.md (the citation check), marked in the text.
\subsection{The compilation: the non-interactive theorems, the numbers and the hidden assumption}\label{sec:gen-opening-analysis}

\begin{sloppypar}
Short forms in this section. Specification files are under \file{doc/leanvm/body/}: \code{03:} is
\file{03-proving-primitives.tex}, \code{08:} is \file{08-end-to-end-protocol.tex}, \code{A:} is
\file{a-ring-switching.tex}, \code{B:} is \file{b-polynomial-commitment-scheme.tex}. Rust files:
\code{whir.rs}, \code{whir_config.rs} are under \file{crates/pcs/src/}; \code{fs/} is
\file{crates/fiat_shamir/src/}; \code{mod.rs} is \file{crates/lean_vm/src/cpu/mod.rs}, and
\code{pcs.rs}, \code{leaf.rs} are under \file{crates/lean_vm/src/}; \code{zerocheck.rs} is
\file{crates/flock/src/zerocheck.rs}. \code{py:} is \file{python-verifier/verifier.py}. These are at
leanVM's pin \code{a386121f}. \code{bp:} is the blueprint at leanerVM \code{b435631};
\code{architecture.md} is \file{docs/architecture.md} there; \code{Compose.lean}, \code{Seams.lean}
and \code{Instance.lean} are under \file{LeanerVM/Protocol/Spine/} at \code{b435631}. ArkLib at its
old pin \code{dca90385} unless a revision is named: \code{Basic.lean}, \code{FiatShamir/Basic.lean},
\code{BCS/Basic.lean}, \code{ProtocolSpec/Basic.lean}, \code{Security/Implications.lean} and
\code{Security/RoundByRound.lean} are under \file{ArkLib/OracleReduction/};
\code{Commitments/Functional/Basic.lean} is under \file{ArkLib/}; \code{CapacityBounds.lean} is
\file{ArkLib/Data/CodingTheory/ProximityGap/CapacityBounds.lean}. The number after the colon is the
line. Other dossiers are cited by file and section: \file{lib-arklib.md} (ArkLib),
\file{literature.md} (the literature), \file{gt-bus.md}, \file{gt-table-pub.md},
\file{gt-flock-ring.md} (the other phases).
\end{sloppypar}

% ---------------------------------------------------------------- E
\subsubsection{Are the blueprint's non-interactive theorems well formed?}\label{sec:gen-opening-wellformed}

The four sketches (\code{bp:1218-1227}, \code{bp:1247-1250} at \code{b435631}):

\begin{leancode}
theorem verify_iff_compiled (prog input proof) :
    verify prog input proof = true ↔
      ∃ s, s.Admissible prog ∧ (Verifier.fiatShamir (leanVmIopp prog s input flock).verifier …)
        accepts (decode s proof) under the BLAKE2s challenge oracle
structure FiatShamirSecurity where …   -- rbr knowledge soundness of the IOPP ⇒ knowledge soundness of its FS compilation in the ROM, error Q · max_i ε_i
structure BcsSecurity where …          -- Merkle-compiled oracles: extraction from collision resistance / ROM
theorem verify_knowledgeSound (fs : FiatShamirSecurity) (bcs : BcsSecurity) (mca : McaJohnson) (flock) :
    … knowledge soundness of `verify` with error niError s Q
theorem baseVerifier_extractsExecution (fs bcs mca flock) (h : verify prog input proof = true) :
    except with probability niError, ∃ t, ValidExecution prog input t
theorem baseProver_complete (hfill : HasFillBlocks prog) (h : ValidExecution prog input t) :
    verify prog input (prove prog input (witness of t)) = true
\end{leancode}
\src{docs/roadmap/protocol-blueprint.md:1218-1227, 1247-1250 at b435631}

They are read here as statements to be written in Lean. Two library objects are introduced first.

\begin{define}{Random oracle; the random-oracle model}
A \textbf{random oracle} is a function drawn uniformly from all functions of a fixed domain and
range, to which every party has query access only. A theorem ``in the random-oracle model'' (ROM)
is a probability statement over that draw, against adversaries that make at most $Q$ queries.
\end{define}

\begin{define}{ArkLib's Fiat--Shamir transform (\lean{Verifier.fiatShamir})}
\lean{Verifier.fiatShamir} (\code{FiatShamir/Basic.lean:129-136} at \code{dca90385}, verbatim in the
probe of \cref{sec:gen-opening-probe-fs}) turns an interactive \lean{Verifier} into a
\lean{NonInteractiveVerifier} whose one message is the tuple of all prover messages, and which
obtains each challenge by querying a \emph{challenge oracle} \code{fsChallengeOracle StmtIn pSpec}
(\code{ProtocolSpec/Basic.lean:849-854}) at the pair (statement, messages so far)
(\lean{challengeOracleInterfaceSR}, \code{:830-834}: ``\code{Query := StmtIn × pSpec.MessagesUpTo i.1.castSucc}'').
That oracle is an \lean{OracleSpec}; any deterministic function of the same signature is a
\lean{QueryImpl} for it, so a concrete hash chain can be plugged in by \lean{simulateQ} (the probe's
\lean{runWithChain}).
\end{define}

% ---------------------------------------------------------------- E.1
\paragraph{\texorpdfstring{\lean{baseVerifier_extractsExecution}}{baseVerifier_extractsExecution}, and what a well-formed statement looks like}\label{sec:gen-opening-extracts}

\paragraph{As sketched, it has no probability space.} \code{verify prog input proof = true} is a
decidable proposition about three fixed values: \lean{verify} is ``a total, computable function of
\code{(prog, input, proof)} to \code{Bool}'' (\code{bp:326}) and BLAKE2s is a fixed function.
``Except with probability \lean{niError}'' then quantifies over nothing: for fixed
\code{(prog, input, proof)} the conclusion \code{∃ t, ValidExecution prog input t} is either true or
false, and no randomness makes the hypothesis probabilistic. The same defect is in the source it
transcribes, the target theorem T4 in \code{architecture.md:264-269} (``\code{BaseVerifier.Accepts version proof program publicInput → except with probability baseError, ∃ assignment, …}''),
whose next sentence asks for what the sketch omits: ``The theorem must expose the component
soundness/knowledge-soundness bounds, Fiat--Shamir or random-oracle model, commitment and hash
assumptions'' (\code{architecture.md:271-273}).

\paragraph{As sketched, it is soundness, not knowledge soundness.} The conclusion is an existential
over traces. Knowledge soundness names an extractor that outputs the witness; the blueprint's own
doctrine (\code{bp:96-103}, acceptance test 24) insists on a named, computable extractor for the
oracle protocol, and then discards it at the last step. Since \lean{ValidExecution} is not what an
extractor can output (\code{bp:96-99} says so), the right shape is: the extractor outputs a stack
$q$; \code{witnessOf q} satisfies \lean{SatisfiedBy} except with the error;
\lean{constraintSoundness} (the target theorem T1) then gives the trace. The knowledge statement is
about the first step; the existential is its corollary.

\paragraph{What the randomness is over, and what the adversary is.} In the ROM the verifier is
$\mathit{verify}^H$, \lean{verify} with the compression function $H : \{0,1\}^{512} → \{0,1\}^{256}$
(the primitive of \code{fs/lib.rs:18-24}, ``\code{f(a, b) = BLAKE2s(a‖b)}'', used for the chain, the
grinding and every Merkle node and leaf, \code{fs/merkle.rs:40-51}) left as an oracle. (The citation
check corrected two points. The quoted text is the doc comment at \code{fs/lib.rs:9};
\code{compress} is at \code{:18-24}. And the one compression does not hash the leaves: a node is one
64-byte BLAKE2s-256, but a leaf is BLAKE2s-256 of the whole leaf image, 512 bytes at level 0 and 384
bytes deeper, eight or six compressions chained through BLAKE2s's own chaining value and counter
(\code{fs/merkle.rs:38-42, 61-67}; \code{primitives/src/hash.rs:209-219}). A random oracle for the
single 64-byte compression does not model the leaf hash; a ROM statement needs either the iterated
construction or a second oracle for leaves.) The adversary
$A^H$ is a $Q$-query oracle algorithm that outputs \code{(prog, input, proof)} (adaptively: leanVM's
statement includes the program, whose hash seeds the chain, \code{mod.rs:82-93}, and the sizes are
the prover's, \code{mod.rs:118-124}). The probability is over $H$ (and $A$'s coins). A well-formed
statement:

\begin{leancode}
theorem verify_knowledgeSound (Q : ℕ) (A : Adversary H Q) :
  Pr_H [ let (prog, input, proof, log) ← A^H ;           -- log: A's queries to H and the answers
         verify^H prog input proof = true ∧
         ¬ SatisfiedBy prog input (witnessOf (extract prog input proof log)) ]
    ≤ niError Q
theorem baseVerifier_extractsExecution (Q) (A) :
  Pr_H [ … verify^H prog input proof = true ∧ ¬ ∃ t, ValidExecution prog input t ] ≤ niError Q
\end{leancode}
\src{the dossier's proposed statements (dossiers/gt-opening-compile.md, section E.1)}

with \lean{extract} a \textbf{straight-line} extractor: it reads the Merkle leaves of the level-0
commitment off the query log (every leaf the prover can open was hashed, so it is in the log except
with the collision or preimage term), list-decodes the word to its $≤ L_0$ nearby codewords
(\cref{sec:gen-opening-hidden}), and returns the member that satisfies \lean{M3Holds}, or a default.
The second theorem follows from the first by \lean{satisfiedBy_witnessOf} and
\lean{constraintSoundness} pointwise. Two further features of the deployed object a faithful
statement carries:
\begin{itemize}
\item \textbf{The sizes are the adversary's.} \lean{verify} dispatches on the eight announced sizes
  (\code{mod.rs:144-149}), so it is the union over admissible $s$ of the members of a protocol
  family; a bound proved per member (per $s$) transfers to the union only by an argument about how
  $A$'s $Q$ queries split among members (they do, since a query to the chain seeded and then fed the
  sizes belongs to one $s$); \lean{niError} must be uniform in $s$ or take the maximum. Acceptance
  test 21 (\code{bp:1326-1328}) records the dispatch and not this obligation.
\item \textbf{The program enters the chain by its hash.}
  \code{iv = BLAKE2s("leanvm" ‖ len ‖ R1CS_DIGEST ‖ bytecode_hash)} (\code{mod.rs:82-93}): two
  programs with the same \code{bytecode_hash} share every challenge, so the statement binds the
  program only up to a BLAKE2s collision, a term of \lean{niError}.
\end{itemize}

\paragraph{The random-oracle heuristic.} \lean{verify} is a function; $\mathit{verify}^H$ is a
family. The only theorem provable is about $\mathit{verify}^H$ with $H$ random; the sentence
``BLAKE2s's compression function is that $H$'' is the \textbf{random-oracle heuristic}, a
non-formal step to be named as such in Layer 13 and in \code{architecture.md} (which already has
the words ``Fiat--Shamir or random-oracle model'' but files them under ``must expose'', not under
``assumed''). It cannot be a Lean hypothesis (it is false for any fixed function, since an
adversary may hard-wire it) and it cannot be an axiom (\file{AGENTS.md}); the theorem about
\lean{verify} itself, with BLAKE2s fixed, is the instantiation \code{H := blake2sCompress} of the
definitions only (\code{verify = verify^blake2sCompress}, by \code{rfl}), never of the security
theorem. The same holds of the Merkle side (collision resistance of BLAKE2s-256 is the second thing
the heuristic buys).

% ---------------------------------------------------------------- E.2
\paragraph{The three assumed interfaces: what each must state to be true and strong enough}\label{sec:gen-opening-interfaces}

\begin{sloppypar}
\paragraph{The ground on ArkLib, at both pins.} The ArkLib dossier (\file{lib-arklib.md},
section G.4) established it; it was re-verified here at \code{dca90385} and \code{7653a901}.
\lean{fiatShamir_completeness} is admitted and stated for the constant challenge oracle
(\code{FiatShamir/Basic.lean:163-173}); soundness after Fiat--Shamir is a \code{TODO} (\code{:175});
round-by-round $⇒$ state-restoration is commented out (\code{Security/Implications.lean:230-254},
and the commented statement has the wrong error, \code{∑ i, rbrError i}, where a state-restoration
bound must depend on the query budget); state-restoration $⇒$ plain knowledge soundness is admitted
(\code{:305-313}); \lean{BCSTransform} is commented out (\code{BCS/Basic.lean:59-63, 79-82});
\lean{Commitment.extractability} has \code{False} for its body
(\code{Commitments/Functional/Basic.lean:250-255}). (The citation check adds the lines at the new pin
\code{7653a901}: \code{Security/Implications.lean:200-224} and \code{:275-283},
\code{Commitments/Functional/Basic.lean:249-254}; the body there is
\code{∃ _extractor, ∀ AuxState _adversary _prover, False}, with the intended probability statement
commented out below it; and the commented round-by-round theorems are about the protocol with added
salts, \lean{addSalt}, which the dossier does not say.) The pull requests the tracker names: \#469
(open since 2026-04-19, ``statements of Section 5'' of [CO25], the duplex-sponge Fiat--Shamir, no
security proof) and \#848 (open, stacked on \#469, ``prove Theorems 6.1 and 6.2'');
\textbf{\#627 is an issue, not a pull request} (a design for the BCS transform, 2026-07-10), whose
staged plan says: ``Full ROM/state-restoration soundness inheritance is deliberately \textbf{out of
scope} and will be fenced as a named open statement''.
\end{sloppypar}

\paragraph{\lean{McaJohnson}.} What it must state: [BCHKS25] Theorem 4.6 as Annex B quotes it
(\code{B:176-185}, \code{thm:mca-johnson}): for the Reed--Solomon code of block length $n$ and
reduced rate $ρ = (k−1)/n$, every radius $δ < 1 − \sqrt{ρ}$, every two-row word $C$, the probability
over $s ← \E$ that $C_s$ agrees with the code on a set of $(1−δ)n$ columns on which $C$ does not
agree with the interleaved code is at most $a/\abs{\E}$, $a$ the explicit expression in
$m = \max(⌈\sqrt{ρ}/(1−\sqrt{ρ}−δ)⌉, 3)$. ArkLib states exactly this, admitted, at both pins:

\begin{leancode}
-- ArkLib/Data/CodingTheory/ProximityGap/CapacityBounds.lean:203-219 (dca90385; 7653a901 has the
-- same statement)
theorem rs_mcaError_le_in_johnson_range
    (domain : ι ↪ F) (k : ℕ) (δ : ℝ≥0)
    (_hk : 1 < k) (_hδ_pos : 0 < δ)
    (_hδ :
        (δ : ℝ) <
          1 - ((((k - 1 : ℕ) : ℝ) / Fintype.card ι) ^ ((1 : ℝ) / 2))) :
    mcaError (AffineLineGenerator F) (ReedSolomon.code domain k) (δ : ℝ) ≤
      ENNReal.ofReal
        (let n : ℝ := Fintype.card ι
         let ρ : ℝ := (k - 1 : ℕ) / n
         let m : ℝ := max ⌈(ρ ^ ((1 : ℝ) / 2)) /
           (1 - ρ ^ ((1 : ℝ) / 2) - δ)⌉ 3
         ((2 * (m + 1/2) ^ 5 + 3 * (m + 1/2) * δ * ρ)
            / (3 * ρ ^ ((3 : ℝ) / 2)) * n
          + (m + 1/2) / ρ ^ ((1 : ℝ) / 2))
           / (Fintype.card F : ℝ)) := by
  sorry -- ABF26-T4.12; external admit [BCHKS25 Thm 4.6].
\end{leancode}
\src{ArkLib at dca90385 (and 7653a901), ArkLib/Data/CodingTheory/ProximityGap/CapacityBounds.lean:203-219}

Instantiated at \code{F := E} and \lean{domain} the additive subspace of $\K ⊂ \E$ (\code{B:90, 302}),
for the affine line $u_0 + s·u_1$, which Annex B's footnote (\code{B:158}) converts to the fold
$(1−s)u_0 + s·u_1$ in characteristic two. It is a statement of mathematics, not about protocols, so
quantifying it universally (over \code{ι}, \code{F}, \lean{domain}, \code{k}, \code{δ}) is right and
true if the theorem is; its consumer is Lemma B.10 (\code{lem:fold-list}; the citation check
corrected the dossier's ``Lemma B.8''), through the $2^{ℓ−1}$ union of rows. Strong enough: yes, it is
the only coding-theoretic input Theorem B.7 uses (\code{B:140-152}). (The citation check qualifies:
the Johnson bound, Theorem B.3, proved at \code{B:391-414}, and Lemma B.11 are also inputs; Theorem
B.9 is the only \emph{external} one.) Its witness obligation is the preprint's theorem (the literature dossier,
\file{literature.md}, section C.3, notes the two readings of the constant $m$; the printed one, the
larger error, is what both Annex B and ArkLib state). Well formed.

\paragraph{\lean{FiatShamirSecurity}.} As the comment states it (``rbr knowledge soundness of the
IOPP $⇒$ knowledge soundness of its FS compilation in the ROM, error $Q · \max_i ε_i$'') it is a
universally quantified claim over IOPPs and is \textbf{false}:
\begin{enumerate}[label=(\alph*)]
\item an IOPP's messages are oracles; its Fiat--Shamir compilation hashes commitments, not oracles,
  and the compiled error carries a hash term (\file{literature.md}, section A.3, item 2:
  $3.5·t^2/2^λ$ in [CY24] Theorem 25.2.1, $3(Q^2+1)/2^κ$ in [BGKTTZ23] Theorem 3.15);
\item leanVM's compilation grinds 17 bits at every query round, which no theorem in the literature
  covers (section A.3, item 3) and which changes the error from $Q·\max ε_i$ to
  $(t + k)·\max_i 2^{−b_i}ε_i$ at best;
\item the hash-chain construction (\code{fs/lib.rs:85-99}, a state ratcheted by $H$) is neither
  ArkLib's per-round oracle keyed by the whole prefix nor [CO25]'s permutation sponge, so a theorem
  about either applies only through a further ``chain is as good as the ideal oracle'' lemma
  (collision term $Q^2/2^{256}$).
\end{enumerate}
\begin{sloppypar}
What ArkLib is about to offer (\#848, \file{SingleSalt.lean:973-993} at its head \code{1c5c5bd7},
\lean{single_salt_fiat_shamir_knowledge_soundness}) takes \textbf{state-restoration} knowledge
soundness of a plain \lean{Verifier} with coins, for the challenge oracle
\code{fsChallengeOracle (StmtIn × Salt) pSpec} implemented by \lean{srChallengeQueryImpl'}, and
concludes adaptive NARG knowledge soundness with the same error;
it has no BCS side, no grinding, and needs round-by-round $⇒$ state-restoration, which is not stated
anywhere. (The citation check corrected the dossier's ``implemented by lazy uniform sampling'':
\lean{srChallengeQueryImpl'} answers from a function that \lean{fsInit} produced ahead of time and
never updates it, \code{ProtocolSpec/Basic.lean:1094-1098} at \code{1c5c5bd7}, ArkLib's docstring
saying ``the whole function is sampled ahead of time''; the distribution is whatever \lean{fsInit}, a
hypothesis, supplies.) To be true and strong enough the interface must be stated for \textbf{this} construction:
(1) a public-coin interactive argument $V^H$ (the Merkle-compiled protocol,
\cref{sec:gen-opening-listbinding}) that is state-restoration knowledge sound against $Q$-query
provers in the ROM with error $ε_{\mathrm{sr}}(Q)$, (2) the hash chain of \code{fs/lib.rs} with $H$
random, (3) the named ground rounds with their bits $b_i$; and conclude adaptive knowledge soundness
of the chain-compiled NARG with error $ε_{\mathrm{sr}}(Q) + c·Q^2/2^{256}$; and the grinding must be
inside $ε_{\mathrm{sr}}$ as ``a query to $H$ at a ground round costs $2^{b_i}$ queries in
expectation''. None of this exists in Lean; the honest interface names the literature and the gap
(\file{literature.md}, section A.3: ``a new assumption to name'').
\end{sloppypar}

\paragraph{\lean{BcsSecurity}.} As the comment states it (``Merkle-compiled oracles: extraction from
collision resistance / ROM'') it names a property with no error and no object. At the pin nothing
can inhabit it, and \#627's plan excludes the ROM inheritance it would need. To be true and strong
enough: for the IOP $C$ of \cref{sec:gen-opening-listbinding} (oracle messages: the level-0 $\K$-word
and the deeper $\E$-words, all with position queries), the Merkle compilation of
\file{crates/pcs/src/merkle.rs} and \code{fs/merkle.rs} (leaf = BLAKE2s of the row image, node =
BLAKE2s of the two children, fixed height and fixed leaf width; the leaf hash being multi-block, the
ROM must model it by a second oracle or through the iterated construction, as the citation check
notes, \cref{sec:gen-opening-extracts}) in the ROM preserves
state-restoration knowledge soundness, with the straight-line extractor that reads every committed
leaf off the query log, up to $3(Q^2+1)/2^{256}$ ([BGKTTZ23] Theorem 3.15's hash term; [CY24]
Theorem 25.2.1). Its hypothesis must include that the verifier fixes the leaf width and the tree
height from public data (checks 12 to 17 of \cref{tab:gen-opening-checks}): that is what makes
leaf/node confusion impossible without domain separation.

\paragraph{Two interfaces the blueprint does not list and that the theorem needs.}
\begin{itemize}
\item \textbf{Round-by-round $⇒$ state-restoration} knowledge soundness, with error $Q_i·ε_i$ summed
  over rounds or $Q·\max_i ε_i$ ([CY24] Theorem 31.2.1/31.3.1: $(t + k)·\max_i ε_i$;
  \file{literature.md}, section A.3, item 1); commented out in ArkLib
  (\code{Implications.lean:230-254}) with the wrong error.
\item \textbf{The list-binding compilation} of \cref{sec:gen-opening-listbinding}, which is not an
  upstream theorem at all.
\end{itemize}

% ---------------------------------------------------------------- E.3
\paragraph{From WHIR's list soundness to extraction of one \texorpdfstring{$q$}{q}}\label{sec:gen-opening-listbinding}

\paragraph{What Theorem B.7 gives.} Round-by-round \textbf{soundness} (Definition B.5 at the pin,
the citation check correcting the dossier's ``Definition B.2''; \code{B:74-84}, an invariant on partial transcripts with an escape probability per verifier message)
for the relation
$R_{\mathrm{open}} = \{(C^{(0)}, \{(W_t, c_t)\}) : ∃ U ∈ L_{γ_0}(C^{(0)}),\ ⟨W_t, g_U⟩ = c_t\ ∀ t\}$
(\code{B:133-136}), where $L_{γ_0}(C^{(0)})$ is the set of interleaved codewords within radius
$γ_0$ of the committed word, of size at most $L_0 = 1/(2η_0\sqrt{ρ_0})$ (\code{B:137}). Its start
invariant is ``every $U ∈ L_{γ_0}(C^{(0)})$ violates at least one of the $J_0$ claims''
(\code{B:218}). No extractor is named: it is soundness for a language, in the IOP model (the oracles
$C^{(i)}$ are read at queried columns only; the prover is unbounded).

\paragraph{What the master theorem gives.} \lean{piop_rbrKnowledgeSoundness}
(\code{Compose.lean:172-177} at \code{b435631}) is ArkLib's worst-case round-by-round knowledge
soundness for \lean{piopExtractor}, whose first step reads the stack $q$ off the commit message
(\lean{commitExtractor}, \code{Compose.lean:58-59}), with a knowledge state function whose
\lean{toFun_full} property says (\code{Security/RoundByRound.lean:186-189} at \code{dca90385}): if
the verifier can output a statement in \code{relOut} with witness \code{witOut}, the state is true at
the full transcript. Everything is relative to \textbf{one} $q$.

\paragraph{The step ``Layer 12 turns list binding into extraction of the one $q$''}
(\code{bp:1190}, acceptance test 11, \code{bp:1296-1299}) is a theorem about the composed IOP, and it
can be stated. Let \lean{Front} be the spine's prefix \code{commit ⟫ bus ⟫ table ⟫ pub ⟫ flock},
whose verifier never queries the oracle (true of the built phases; a requirement on the others under
the inner-product interface of \cref{sec:gen-opening-alternatives}), with worst-case round-by-round
knowledge soundness from \lean{M3Rel} to \lean{Seam.flock} at errors $ε_i$ (the composition of the
first four fields of \lean{Phases.Security}). Let \lean{Front'} be \lean{Front} with the commit
message's type replaced by an interleaved $\K$-word $w$ (the level-0 codeword, with position
queries) and the statement carried unchanged; since the front's verifier never reads the oracle,
\lean{Front'}'s verifier is \lean{Front}'s. Let \code{C := Front' ⟫ whirOpen} (the level-0 batching
of \lean{whirOpen} being the opening phase's $λ$). Then:

\begin{define}{List-binding compilation (the theorem the dossier proposes)}
$C$ is worst-case round-by-round knowledge sound for \lean{M3Rel} with extractor
$\mathcal{E}(w, τ) :=$ the first $g ∈ Λ(w)$ with \code{M3Holds input g} (else a default), where
$Λ(w) = \{g_U : U ∈ L_{γ_0}(w)\}$, and per-challenge errors $L_0·ε_i$ at the front's challenges and
Theorem B.7's at WHIR's. Its knowledge state function is: after $w$, ``no $g ∈ Λ(w)$ satisfies
\lean{M3Holds}''; inside the front, ``for every $g ∈ Λ(w)$, the front's state at $(g, τ)$ is false''
(the front's state function evaluated with $g$ in place of the commit message); at the seam this
reads ``for every $g ∈ Λ(w)$ some pooled claim is false of $g$'', which is $∉ R_{\mathrm{open}}$ and
Theorem B.7's start invariant; then Theorem B.7's invariant.
\end{define}

\paragraph{Proof shape.} The front's \lean{toFun_next} and \lean{toFun_full} transfer pointwise
(each $g$ is a transcript of \lean{Front}); the escape probability of a front challenge is bounded by
the union over $Λ(w)$ of the front's worst-case bound, which is per transcript and therefore per $g$
(\lean{rbrKnowledgeSoundnessWorstCaseWith}, \code{Security/RoundByRound.lean:553-568}: a bound for
every statement, every round, every partial transcript); the seam step needs \lean{toFun_full} of
the front's state function in the contrapositive: state false at the front's last transcript $⇒$ the
verifier's output claims are not all true of $g$. That is exactly why the spine's \emph{worst-case}
form is the right one (the averaged \lean{rbrKnowledgeSoundness} would not union-bound). What it
needs and does not have: Theorem B.7 in Lean (the hole for it, K1); the lemma that every member of
$Λ(w)$ is $\K$-valued and of the right size (\cref{sec:gen-opening-hidden}); a generic ``state
functions indexed by a finite set'' lemma (union bound; ArkLib has none); and the identification of
the front's oracle-free verifier with \lean{Front'}'s (a cast of \lean{OracleReduction}, which ArkLib
\#383's \file{Cast.lean} is about). The extractor is computable as a definition (enumerate the
codewords of the finite code; Lean has no cost model, \code{bp:102-103}), and efficient in the
literature's sense through Guruswami--Sudan, which Lean need not implement.

This theorem is where the factor $L_0$ enters \lean{niError} (\cref{sec:gen-opening-numbers}), and
it is the missing fourth ingredient of \lean{verify_knowledgeSound}: neither
\lean{FiatShamirSecurity}, \lean{BcsSecurity} nor \lean{McaJohnson} contains it, and the sentence
``\lean{verify_knowledgeSound} extracts the member'' (\code{bp:1299}) assigns it to a theorem whose
sketched hypotheses cannot prove it.

% ---------------------------------------------------------------- E.4
\paragraph{\texorpdfstring{\lean{verify_iff_compiled}}{verify_iff_compiled}}\label{sec:gen-opening-iff-compiled}

``Unconditional: it is about two definitions'' (\code{bp:1231-1232}) is right in kind and wrong in
the definitions named.
\begin{sloppypar}
\begin{enumerate}
\item \textbf{\lean{Verifier.fiatShamir} takes a \lean{Verifier}, not an \lean{OracleVerifier}}, and
  an oracle verifier becomes one only through \lean{OracleVerifier.toVerifier}
  (\code{Basic.lean:490-496} at \code{dca90385}; the probe of \cref{sec:gen-opening-probe-fs} prints
  its type), whose statement is \code{StmtIn × (∀ i, OStmtIn i)} and whose transcript holds the
  oracle messages \textbf{in the clear}: \code{(leanVmIopp …).verifier.toVerifier} has the level-0
  word ($2^{μ−6}·2^6$ elements of $\K$) and every deeper codeword as messages, and
  \lean{fsChallengeOracle} would hash them whole. No \code{decode : Proof → messages} exists: a proof
  holds roots and opened rows (\code{fs/transcript.rs:9-12}), from which no codeword is recoverable.
  The right-hand side must be the Fiat--Shamir compilation of the \textbf{Merkle-compiled}
  interactive verifier, whose messages are the roots, the scalars, the grinding nonces and the
  openings; ArkLib has no such compilation (\lean{BCSTransform} commented out), so Layer 12 must
  define it (\lean{bcsCompile}) before it can state the theorem.
\item \textbf{The challenge oracle can be the chain.} The probe of \cref{sec:gen-opening-probe-fs}
  shows the mechanism: a function \code{chain : (i) → StmtIn × MessagesUpTo i → Challenge i} is a
  \code{QueryImpl (fsChallengeOracle …) Id}, and \lean{simulateQ} runs the compiled verifier under it
  to an \code{Option StmtOut}. leanVM's chain is such a function once the statement is
  \code{(input)} with \code{(prog, s)} as parameters of the family (the chain's seed depends on
  \code{prog} through \code{iv}, \code{mod.rs:82-93}, and its first eight absorbs are the sizes,
  \code{mod.rs:118-124}) and the messages are absorbed in the wire encoding (one coefficient of each
  round polynomial dropped, \code{fs/transcript.rs:63-71}, \code{:289-309}; a root as two scalars,
  \code{fs/merkle.rs:14-20}; the Merkle openings \textbf{not absorbed},
  \code{fs/transcript.rs:223-224} ``Not absorbed: its binding is the Merkle structure itself'', the
  citation check correcting the dossier's \code{:225-228}). A
  challenge that is several squeezes ($α ∈ \E^4$, the query positions by chunking,
  \code{whir.rs:1175-1193}) is one \code{Challenge i} of a product type, or several rounds.
\item \textbf{The grinding checks are not in the right-hand side.} The nonce at each query round is
  a prover message (absorbed with tag 4, \code{fs/lib.rs:118-120}), and the verifier's check
  \code{pow_bits_ok(pow_base(cv), nonce, 17)} (\code{fs/lib.rs:47-51, 165-174}) reads the
  \textbf{chain state} \code{cv}, which an ArkLib \lean{Verifier} does not see: its \lean{verify}
  takes the statement and the full transcript (\code{Basic.lean:248-250}). Either the interactive
  verifier recomputes the chain from the transcript (definable, since the chain is a function of
  statement and messages; it then depends on $H$ and is no longer an IOP verifier), or
  \lean{verify_iff_compiled}'s right-hand side is a conjunction ``the compiled verifier accepts
  $∧$ every grinding check passes'', the second conjunct being about the chain. The sketch has
  neither.
\item \textbf{\code{∃ s, s.Admissible prog}} is right in shape (the family is indexed by the sizes;
  the proof carries them as its first eight scalars) provided \lean{Admissible} is the verifier's
  whole predicate, the stacking window and the rate window included (\cref{f:bus-admissible}).
\item \textbf{The equation is a definition, not a theorem, if \lean{verify} is defined as the
  right-hand side.} The blueprint wants \lean{verify} ``written from §8.5 before
  \code{cpu/mod.rs:711-779} is opened'' (\code{bp:1216}) and the theorem to ``pin the shape''
  (\code{bp:1277}, \code{:1302}, \code{:1316}, \code{:1328}). That is only a check if the two sides
  are written independently; then the theorem is provable by unfolding (\code{rfl}-like), and its
  value is the review of the two definitions.
\end{enumerate}
\end{sloppypar}

Statable form (with the inner-product interface and the corrections above):

\begin{leancode}
theorem verify_iff_compiled (prog input proof) :
    verify prog input proof = true ↔
      ∃ s, s.Admissible prog ∧ ∃ msgs, decode prog s proof = some msgs ∧
        runWithChain (bcsCompile (leanVmIopp prog s)).verifier (blake2sChain prog s) input msgs
          = some () ∧ grindingChecks prog s msgs = true
\end{leancode}
\src{the dossier's proposed statement (dossiers/gt-opening-compile.md, section E.4)}

% ---------------------------------------------------------------- E.5
\paragraph{Completeness}\label{sec:gen-opening-completeness}

``Perfect completeness survives the transform without any assumption'' (\code{bp:1256-1257}) is
false for the deployed transform, for two reasons of different weight.

\paragraph{Grinding.} The honest prover's grind is a search: \code{fs/lib.rs:126-155}
\code{grind_pow} loops over nonces $0, 1, 2, …$ ``until'' \code{pow_bits_ok}
(\code{loop { … break n }} at \code{:132-138}, and the parallel form at \code{:139-151}), with no
bound. For a fixed state \code{cv} there is no theorem that a \code{u64} nonce (or any nonce) with
\code{compress(pow_base(cv), nonce ‖ 4)[0] ≡ 0 mod 2^17} exists; in the ROM the search fails with
probability $(1 − 2^{−17})^{2^{64}}$, not zero; with BLAKE2s fixed the statement ``for every reachable
\code{cv} a nonce exists'' is a conjecture about BLAKE2s. So \code{prove prog input w} in Lean is
either partial (\code{Option Proof}, returning \code{none} when a fuelled search fails) or takes the
nonces as an argument, and \lean{baseProver_complete} is
\code{prove … = some proof → verify … proof = true} (or ``for every choice of nonces that passes'').
``Without any assumption'' must go; the assumption is the success of the search, and it is harmless
only because it is checked by running the prover.

\paragraph{Phases complete up to an error.} \lean{Phases.Complete} demands
\lean{perfectCompleteness} of every phase (\code{Compose.lean:94-104};
\lean{Component.Complete.complete} is \lean{perfectCompleteness}, \file{lib-arklib.md}, section G.2).
If any phase were complete only up to an error, \lean{Phases.Complete} would be uninhabitable and
\lean{baseProver_complete}, an equation, unprovable. The blueprint's acceptance test 20
(\code{bp:1323-1325}) says Flock has an exceptional challenge; the dossier on Flock settles it
(\file{gt-flock-ring.md}, section 4; \cref{f:flock-ring-test20-completeness}): the \emph{protocol}
is perfectly complete (no verifier takes an inverse; the derived coefficient
$c_0 = \mathit{claim} + r·(c_1 + …)$ is right at every $r$, \code{fs/transcript.rs:297-302}), and the
exceptional point belongs to the pinned Rust prover's shortcut (\code{zerocheck.rs:117}). So the
Lean honest prover can be perfectly complete phase by phase, and the composition
(\lean{Reduction.seqCompose_perfectCompleteness_of_pure}, \code{bp:242}) applies; the ArkLib
dossier's proposal of an error field in \lean{Component.Complete} (\file{lib-arklib.md}, section
G.2) remains the right insurance, and with it ``survives the transform'' becomes ``with the sum of
the phases' errors, plus the grind''.

\paragraph{What does survive.} With the chain deterministic, an honest proof that the interactive
verifier accepts on the chain's challenges is accepted by \lean{verify}: that is a statement about
two definitions (the interactive run with challenges supplied by the chain equals the compiled
run), true by unfolding, and the grinding is the only randomness left in the honest prover.

% ---------------------------------------------------------------- F
\subsubsection{The numbers}\label{sec:gen-opening-numbers}

\paragraph{The blueprint's \texorpdfstring{\lean{piopError_le}}{piopError_le} and acceptance test 23}\label{sec:gen-opening-piop-error}

\begin{sloppypar}
\code{piopError_le (hs : s.Admissible prog) : Σ i, piopError s i ≤ 2^40/|E| + flockError}
(\code{bp:1130}); test 23: ``at the caps, $Σ\,\mathit{piopError} < 2^{-150}$ plus Flock's
$< 2^{-183}$'' (\code{bp:1331-1332}). Evaluated here with the blueprint's own per-phase charges, at
the largest admissible instance. The sizes are bounded by the stacking window,
$μ_{\mathrm{stack}} ≤ 28$ (\code{mod.rs:174-176}, \code{pcs.rs:51}): every column has at most
$2^{28}$ entries, so $κ_{\mathrm{mem}} ≤ 26$, $τ_j ≤ 24$, and the bus depth is $μ_{\mathrm{bus}} ≤ 28$
(the Python layout code run over a covering grid of sizes finds 28 as the maximum, at
$κ_{\mathrm{mem}} = 16$, $τ = (0, 0, 23, 0, 23, 3)$, $κ_{\mathrm{bc}} = 26$, scratch
\file{bus_depth.py}, \cref{sec:gen-opening-probe-params}; the bus dossier, \file{gt-bus.md},
section D.4, proves $μ_{\mathrm{bus}} ≤ 28$ by counting). The table sumcheck has $B = 2$ constraints
in all (\code{py:836-843}, the citation check correcting the dossier's \code{py:836-851}: only
\code{JUMP} has constraints, \code{n_constraints = 2};
\code{xi_powers = powers(xi, n_constraints + 3)}, \code{py:1391}), $τ_{\max} ≤ 24$, and $J = 113$.
\end{sloppypar}

{\small\setlength{\tabcolsep}{4pt}
\begin{longtable}{@{}>{\raggedright\arraybackslash}p{0.47\linewidth}>{\raggedright\arraybackslash}p{0.15\linewidth}>{\raggedright\arraybackslash}p{0.32\linewidth}@{}}
\caption{The blueprint's phase errors at the largest admissible instance.}\label{tab:gen-opening-phase-errors}\\
\toprule
Phase, charge (blueprint) & Where & Value at the maximum, as a multiple of $1/\abs{\E}$ \\
\midrule
\endfirsthead
\toprule
Phase, charge (blueprint) & Where & Value at the maximum, as a multiple of $1/\abs{\E}$ \\
\midrule
\endhead
\bottomrule
\endfoot
bus, $4·2^{μ_{\mathrm{bus}}}$ on $(α, β)$ & \code{bp:974} & $2^{30}$ \\
bus, GKR: $(n_{\mathrm{side}} − 1)$ per combiner, $5$ per round, $2$ per pair & \code{bp:941-942}
  & $≈ 14$ layers $× 14 ≈ 2^{7.6}$ \\
zerocheck escape, $τ_{\max}$ per constraint & \code{bp:1001-1003} & $48$ \\
table, $(B + 2)$ on $ξ$, $3$ per round & \code{bp:994} & $4 + 72$ \\
public input, $1$ & \code{bp:1027} & $1$ \\
Flock, $(4·k_{\mathrm{batch}} + 163) + 2^{32}$ & \code{bp:1085} & $≈ 2^{32}$ \\
opening (Layer 10 as sketched), $(J − 1)$ on $λ$, $2$ per round & \code{bp:1114} & $112 + 56$ \\
\midrule
\textbf{sum} & & $≈ 1.25·2^{32} ≈ 2^{32.3}$, that is $≈ 2^{-159.7}$ \\
\end{longtable}
}

The citation check recounts the GKR row: with radix-4 layers of $0, 2, …, 26$ variables the layer
sumchecks have $Σ_{j<14} 2j = 182$ rounds, so $5·182 + 2·14 + 2·14 ≈ 966 ≈ 2^{9.9}$, not $≈ 196$ (its
own reading of \code{bp:940-942}, not checked against a GKR implementation). Immaterial: the sum
stays $≈ 1.25·2^{32}$.

So under the window \lean{piopError_le} holds with a margin of $2^8$ (the citation check: that
compares the whole sum with $2^{40}$; the bound's non-Flock part, $≈ 2^{30}$, sits $2^{10}$ under
$2^{40}$), and the interactive error of
the oracle protocol is about $2^{-160}$, dominated by ring switching's $2^{32}/\abs{\E}$ (which test
23 leaves out of ``Flock's $< 2^{-183}$'': \cref{f:flock-ring-test23-ring-switching}). At the
per-log caps without the window ($κ_{\mathrm{mem}} = τ_j = 32$), $μ_{\mathrm{bus}} = 38$ (scratch run,
\code{stack_log 41 bus depth 38}) and the $(α, β)$ term alone is $4·2^{38} = 2^{40}$, so the bound is
false there, as the bus dossier found (\file{gt-bus.md}, section D.4; \cref{f:bus-admissible}). The
statement is only true if \lean{Admissible} contains the window.

\paragraph{What it measures.} $Σ\,\mathit{piopError}$ is the \emph{interactive} error of the
\emph{oracle} protocol (the sum over challenges, \code{bp:318}). It is not the quantity that governs
the deployed security: after Fiat--Shamir the governing quantity is the largest per-challenge error
(\code{B:84}, \code{whir_config.rs:25-27}), and the deployed challenges with the largest errors are
WHIR's, which the oracle protocol does not contain. $2^{-150}$ is therefore a fact about a
sub-protocol, of no consequence for the 128-bit claim.

\paragraph{What the deployed parameters claim, and in which regime}\label{sec:gen-opening-deployed-claim}

\begin{sloppypar}
\code{SECURITY_BITS = 128} (\code{whir_config.rs:38}; \file{crates/lean_vm/src/lib.rs:69}) is a
\textbf{round-by-round} target: ``every verifier-challenge transition must have conditional failure
probability at most \code{2^-SECURITY_BITS}'' (\code{whir_config.rs:36-37}), ``the Fiat--Shamir error
per random-oracle query is the MAX of the entries, not their sum'' (\code{:25-27}), ``not a claim that
the sum of all interactive failure probabilities is bounded'' (\code{:404-407}). The regime is
\textbf{Johnson list decoding with an explicit slack}, provable given [BCHKS25] Theorem 4.6
(\code{:314-337}: ``That analysis is always the Johnson radius with explicit slack \code{eta} … WITH
out-of-domain binding''; \code{analysis_version}
\code{"bchks25_thm_4_6_exact_reduced_rate_row_union_optimized_eta"}, \code{:961}); no capacity
conjecture, no unique-decoding fallback in production (\code{udr_queries} is
\texttt{\#[cfg(test)]}, \code{:175-180}).
\end{sloppypar}

The per-level parameters are the output of a floating-point search (\code{optimize_johnson_level},
\code{:639-707}) over the theorem parameter $m ≥ 3$, minimizing the query count subject to four
per-level bounds; the Python tabulates the resulting query counts (\code{py:910}). A scratch port of
the search (\file{whir_params.py}, \cref{sec:gen-opening-probe-params}) \textbf{reproduces the Python
table exactly} for all $μ ∈ [15, 28]$ and $ρ ∈ \{1, …, 4\}$ (the ladder geometry, the rates and the
query counts), so the blueprint's \code{ladder_queries_eq : (ladder 15 1).map (·.queries) = [223, 55]}
(\code{bp:1165}) is the right pin and the Rust and the Python agree. What the search yields, at three
sizes:

\begin{landscape}
{\footnotesize\setlength{\tabcolsep}{3.5pt}
\begin{longtable}{@{}llllrllrrrrrrr@{}}
\caption{The deployed WHIR parameters at three sizes: the analysis parameters, the list bound and the
security bits of each term.}\label{tab:gen-opening-params}\\
\toprule
$μ$, $ρ$ & level & fold & rate & $m$ & $η$ & $γ$ & $L ≤ 1/(2η\sqrt{ρ})$ & queries & OOD & fold (MCA) bits
  & query bits ($+17$ grind) & OOD bits & list-unioned algebraic bits \\
\midrule
\endfirsthead
\toprule
$μ$, $ρ$ & level & fold & rate & $m$ & $η$ & $γ$ & $L ≤ 1/(2η\sqrt{ρ})$ & queries & OOD & fold (MCA) bits
  & query bits ($+17$ grind) & OOD bits & list-unioned algebraic bits \\
\midrule
\endhead
\bottomrule
\endfoot
15, 1/2 & 0 & 6 & 1/2 & 395 & 0.00179 & 0.292 & 396 & 223 & 0 & 132.9 & 111.0 & 179.5 & 152.4 \\
        & 1 & 4 & 1/16 & 306 & 0.00080 & 0.753 & 2527 & 55 & 1 & 133.2 & 111.0 & 167.2 & 149.7 \\
\midrule
28, 1/2 & 0 & 6 & 1/2 & 110 & 0.00643 & 0.286 & 110 & 228 & 0 & 129.2 & 111.0 & 180.4 & 154.2 \\
        & 1 & 4 & 1/16 & 81 & 0.00309 & 0.747 & 648 & 56 & 1 & 129.8 & 111.0 & 169.9 & 151.7 \\
        & 2 & 4 & 1/128 & 46 & 0.00192 & 0.910 & 2944 & 32 & 1 & 130.4 & 111.0 & 165.8 & 149.5 \\
        & 3 & 4 & 1/1024 & 8 & 0.00390 & 0.965 & 4100 & 23 & 1 & 139.2 & 111.1 & 165.2 & 149.0 \\
        & 4 & 4 & $2^{-13}$ & 4 & 0.00274 & 0.986 & 16644 & 18 & 1 & 140.2 & 111.4 & 161.6 & 147.0 \\
        & 5 & 4 & $2^{-16}$ & 5 & 0.00068 & 0.996 & 218453 & 14 & 1 & 134.7 & 111.2 & 154.9 & 143.3 \\
\midrule
28, 1/16 & 0 & 6 & 1/16 & 27 & 0.00926 & 0.741 & 216 & 57 & 0 & 131.7 & 111.0 & 179.4 & 153.3 \\
\end{longtable}
}
\end{landscape}

\begin{sloppypar}
The residual is 5 variables at $μ = 15$ and 2 at $μ = 28$; ``bits'' is $−\log_2$ of the bound; the
columns are \code{whir_config.rs}'s \code{paper_johnson_log_a} (\code{:484-490}, Theorem 4.6's $a$
times the row union $2^{ℓ−1}$ of Lemma B.10, the citation check correcting the dossier's ``Lemma
B.8''), \code{paper_per_query_bits × queries} (\code{:494-498}),
\code{paper_ood_bits} (\code{:600-608}) and \code{johnson_algebraic_bits} (\code{:556-568}).
\end{sloppypar}

Reading:
\begin{enumerate}
\item the \textbf{query rounds} are the weakest: $(1 − γ)^t ≈ 2^{-111}$ at every level, by design
  (\code{:57-60}: ``the query count only needs to close the remaining \code{SECURITY_BITS − 17}
  bits''), the other 17 bits being the proof of work ground before the positions are sampled;
\item the \textbf{fold challenges} sit at $128.2$ to $146$ bits in the sample printed (the
  mutual-correlated-agreement term, with the $2^{ℓ−1}$ row union). (The citation check: over all 56
  supported sizes they run from 128.19, at $μ = 27$, $ρ = 1/8$, level 0, to 148.68, at $μ = 17$,
  $ρ = 1/4$, level 2. And it corrected the dossier's ``the search stops at the first $m$ that clears
  128, \code{:664-665}'': the loop at \code{:662-666} breaks at the first $m$ whose
  mutual-correlated-agreement term \emph{fails} to clear 128, and among the $m$ tried it keeps the one
  with the fewest queries, \code{:697}, which is the largest $m$ still clearing 128; that is why the
  fold bits sit just above 128.)
\item the Johnson list bound is \textbf{$L_0 ≤ 396$ at level 0} ($2^{8.6}$, at $μ = 15$) down to
  $110$ at $μ = 28$, at rate $1/2$, and grows to $2^{17.7}$ at the deepest levels. (The citation check:
  over all rates $L_0$ reaches $648$, at $μ = 20$, $ρ = 1/16$, that is $2^{9.34}$; so $L_0$ runs from
  110 to 648 over the supported sizes.)
\item the algebraic checks ``unioned over the list'' (\code{:537-555}: the batch polynomial's degree
  and the ring switching's $2^{31} + …$, times $L_i$) sit at $143$ to $154$ bits (the citation check:
  $141.0$ to $154.2$ over all sizes, the dossier's own output having $141.00$ and $141.42$ at
  $μ = 28$, level 5); this is the \textbf{only} place the Rust multiplies an oracle-protocol error by
  the list size. The bus's $4·2^{μ_{\mathrm{bus}}}$ is not unioned (\code{leaf.rs:108-118},
  \code{soundness_degree_bound} has no $L$): with $L_0 ≤ 2^{8.6}$ and $μ_{\mathrm{bus}} ≤ 28$ it would
  still clear 128 ($192 − 30 − 8.6 = 153$), so the omission is not a vulnerability at these
  parameters, but it is an omission of the same kind the blueprint makes
  (\cref{f:opening-rust-list-accounting}). (The citation check: the Rust's bound is
  $(N_{\mathrm{TUPLE\_BITS}} + 1)·2^μ + 8(μ+1)^2 = 5·2^μ + …$, the ``$4·2^{μ_{\mathrm{bus}}}$'' being
  the blueprint's term; with $L_0 ≤ 648$ and $5·2^{28}$, $192 − 30.3 − 9.3 ≈ 152$, and the conclusion
  stands.)
\end{enumerate}

\paragraph{\texorpdfstring{\lean{niError}}{niError} against this}\label{sec:gen-opening-nierror}

\lean{niError} ``is the round-by-round maximum times the query bound plus the grinding-adjusted WHIR
terms'' (\code{bp:1234-1235}). Set against \cref{sec:gen-opening-deployed-claim} and the literature
dossier (\file{literature.md}, section A.3):
\begin{itemize}
\item the ``round-by-round maximum'' of the oracle protocol is $≈ 2^{-160}$
  (\cref{sec:gen-opening-piop-error}); of the compiled protocol it is $2^{-111}$ at every query round
  before grinding and $2^{-128.2}$ at the worst fold. ``Plus the grinding-adjusted WHIR terms''
  acknowledges that the maximum is elsewhere but gives no formula; the phrase ``grinding-adjusted''
  has no theorem behind it (section A.3, item 3), and the adjustment is multiplicative on a per-round
  error ($2^{-17}·2^{-111}$), not additive;
\item the factor $L_0$ on every oracle-protocol challenge (\cref{sec:gen-opening-listbinding,sec:gen-opening-hidden})
  is absent; numerically benign ($2^{-160+8.6}$; with the citation check's $L_0 ≤ 648$,
  $2^{-160+9.3}$), structurally load-bearing;
\item the hash term ($Q^2/2^{256}$-type) is absent; numerically it is what caps the 128 bits at a
  256-bit digest (section A.3, item 2);
\item the prover's choice of sizes (\cref{sec:gen-opening-extracts}) is absent.
\end{itemize}

\paragraph{Does the blueprint's bound mean something at the deployed parameters?}
$Σ\,\mathit{piopError} < 2^{-150}$ is a true statement about the ideal-oracle protocol at admissible
sizes and says nothing about \lean{verify}'s security level, which is set by the query rounds at 111
bits of information-theoretic soundness plus 17 bits of work, in the Johnson regime, conditional on
Theorem 4.6. A \lean{niError} that could carry leanVM's claim would read, per random-oracle query,
\[
  \max\Bigl( L_0·\max_{\mathrm{front}} ε_i,\ \max_{\mathrm{fold}}\bigl(2L_i/\abs{\E} + 2^{ℓ_i−j}a_i/\abs{\E}\bigr),\
  \max_{\mathrm{query}} 2^{-b}(1−γ_i)^{t_i},\ … \Bigr)
\]
times $(Q + \mathit{rounds})$, plus the hash term, and it would be a statement about the ROM
construction of \cref{sec:gen-opening-interfaces}, not about the oracle protocol.

% ---------------------------------------------------------------- G
\subsubsection{The oracle model's hidden assumption}\label{sec:gen-opening-hidden}

In the oracle protocol the prover's first message is \code{q : Column μ}, a table of $2^μ$ values
of $\K$, and the verifier reads it only through the interface (evaluation queries today, weighted
queries under the inner-product interface of \cref{sec:gen-opening-alternatives}); the extractor
reads it whole (\lean{commitExtractor}). The compiled protocol replaces $q$ by a Merkle root. For
the compiled protocol to inherit the oracle protocol's security, the commitment must provide:
\begin{enumerate}
\item \textbf{binding to a $\K$-valued table of size $2^μ$}, the object \lean{M3Holds} is about
  (\code{M3Rel I} is on \code{Column I.μ}, \code{Instance.lean:230} at \code{b435631}; the bus's
  fingerprints, the counts, the packing of $\qflock$ all use that the entries are in $\K$);
\item \textbf{weighted openings over $\E$}: the pooled claims have weights and values in $\E$
  (\code{03:113});
\item an \textbf{extractor} that recovers the table from what the adversary did, straight-line.
\end{enumerate}

What the deployed WHIR provides:
\begin{itemize}
\item \textbf{A $\K$-valued word, not a codeword.} The root commits (up to collisions) to a leaf
  image per position: $2^6$ words of 8 bytes at level 0 (\code{whir.rs:331-345},
  \file{crates/pcs/src/merkle.rs:144-192}), and every 8-byte pattern is an element of $\K$
  ($\K = \F_2[x]/(x^{64} + …)$, \code{A:8}; \code{F64(u64)}), so the committed object is an
  interleaved word $w ∈ \K^{2^6 × n_0}$ by encoding, with no check that any row is a codeword:
  \textbf{proximity, not membership} (\code{B:4} ``our PCS is only list binding''; \code{B:70-72}
  Definition B.4, the citation check correcting the dossier's ``Definition B.1''). The verifier supplies the zero prefix of the absent lanes
  (\code{fs/merkle.rs:53-58}, \code{whir.rs:2125-2132} with \code{row_words = n_lanes}), so $w$'s rows
  past $n_{\mathrm{lanes}}$ are zero in the Rust's proof format (\cref{sec:gen-opening-proof} for the
  Python's).
\item \textbf{A list, over $\E$.} Theorem B.7 binds the prover to
  $Λ(w) = \{g_U : U ∈ L_{γ_0}(w)\}$, the interleaved codewords of
  $\mathrm{RS}[\E, \mathrm{dom}_0, 2^{κ_0}]$ within radius $γ_0$ (\code{B:137}), of size at most $L_0$
  ($110$ to $648$ over the supported sizes; the citation check corrected the dossier's ``110 to 396'',
  the range at rate $1/2$ only, \cref{sec:gen-opening-deployed-claim}). A priori these are $\E$-valued multilinears; the
  relation needs $\K$-valued ones. \textbf{They are $\K$-valued}, by the lemma below (the dossier's
  argument; the specification does not say it). So every member of $Λ(w)$ is $\K$-valued, of $μ$
  variables after the relayout $f(u, x) = q(x, u)$ (\code{B:380}), with zero rows where $w$ has them.
  This lemma belongs to Layer 11 or 12 and is a hypothesis of the list-binding compilation
  (\cref{sec:gen-opening-listbinding}); without it the extracted witness has the wrong type.
\item \textbf{Openings over $\E$ of weighted claims}: provided, that is Theorem B.7's relation.
\item \textbf{Extraction}: from the query log (the leaves) and list decoding, as in
  \cref{sec:gen-opening-listbinding}; in the Johnson regime the list is small and Guruswami--Sudan
  finds it; in Lean, enumeration suffices.
\end{itemize}

\begin{define}{Lemma: every codeword close to a $\K$-valued word is $\K$-valued (the dossier's two-line argument)}
Write $\E = \K ⊕ \K·y ⊕ \K·y^2$ and a row $U_u = U^{(0)}_u + y·U^{(1)}_u + y^2·U^{(2)}_u$. The
encoder's generator matrix is $\K$-valued (the novel basis evaluated at points of $\mathrm{dom} ⊂ \K$,
\code{B:302-306}, ``a $\K$-valued message gives a $\K$-valued word''), so each coordinate projection
$U^{(j)}_u$ is a codeword of $\mathrm{RS}[\K, \mathrm{dom}_0, 2^{κ_0}]$. On the agreement set $A$,
$\abs{A} ≥ (1−γ_0)n_0 > ρ_0 n_0 = 2^{κ_0}$ (since $γ_0 < 1 − \sqrt{ρ_0} ≤ 1 − ρ_0$), $U_u = w_u ∈ \K$,
so $U^{(1)}_u$ and $U^{(2)}_u$ vanish on more than $2^{κ_0} − 1$ points and are zero; $U = U^{(0)}$.

The citation check re-derived the argument and finds it sound as written.
\end{define}

\paragraph{The assumption, stated.} The master theorems assume the verifier's oracle \emph{is} one
table $q$. The commitment gives a set $Λ(w)$ of at most $L_0$ tables, fixed when the root is sent,
among which the prover may choose after seeing challenges. The security of the compiled protocol is
the master theorem applied to every member at once, with a union bound: the per-challenge errors of
every phase before the opening are multiplied by $L_0$, the extractor is ``list-decode, then select
the member that satisfies \lean{M3Holds}'', and the opening's job is to show that some member
satisfies every pooled claim (Theorem B.7's relation $R_{\mathrm{open}}$, \code{B:133-136}; acceptance
test 11, \code{bp:1296-1299}, has the right words). The Rust's own accounting states the same in one
place (``the opening's own evaluation claim sits at a post-commit random point, so at most one list
member matches it except with \code{L·μ/|F|}'', \code{whir_config.rs:593-599}) and Flock in another
(\file{literature.md}, section A.3, item 4, [BRW26] Remark 11). Neither the specification (§8.4 is
\code{TODO}, \code{08:44-46}), nor the blueprint's \lean{verify_knowledgeSound}, nor any of its three
interfaces contains it; the sentence ``Layer 12 turns list binding into extraction of the one $q$
that satisfies every pooled claim'' (\code{bp:1190}) names the step without the theorem
(\cref{f:opening-list-binding}).
